\documentclass[10pt,a4paper]{amsart}
\usepackage{amsmath,amssymb,amsthm,graphicx}
\usepackage{verbatim}
\newenvironment{reference}{\comment}{\endcomment}
%\newenvironment{reference}{}{}
\newenvironment{slogan}{\comment}{\endcomment}
\newenvironment{history}{\comment}{\endcomment}

% For commutative diagrams we use Xy-pic
\usepackage[all]{xy}

% We use 2cell for 2-commutative diagrams.
\xyoption{2cell}
% Theorem environments.
%
\theoremstyle{plain}
\newtheorem{theorem}[subsection]{Theorem}
\newtheorem{proposition}[subsection]{Proposition}
\newtheorem{lemma}[subsection]{Lemma}

\theoremstyle{definition}
\newtheorem{definition}[subsection]{Definition}
\newtheorem{example}[subsection]{Example}
\newtheorem{exercise}[subsection]{Exercise}
\newtheorem{situation}[subsection]{Situation}

\theoremstyle{remark}
\newtheorem{remark}[subsection]{Remark}
\newtheorem{remarks}[subsection]{Remarks}

\numberwithin{equation}{subsection}

% Macros
%
\def\lim{\mathop{\mathrm{lim}}\nolimits}
\def\colim{\mathop{\mathrm{colim}}\nolimits}
\def\Spec{\mathop{\mathrm{Spec}}}
\def\Hom{\mathop{\mathrm{Hom}}\nolimits}
\def\Ext{\mathop{\mathrm{Ext}}\nolimits}
\def\SheafHom{\mathop{\mathcal{H}\!\mathit{om}}\nolimits}
\def\SheafExt{\mathop{\mathcal{E}\!\mathit{xt}}\nolimits}
\def\Sch{\mathit{Sch}}
\def\Mor{\mathop{\mathrm{Mor}}\nolimits}
\def\Ob{\mathop{\mathrm{Ob}}\nolimits}
\def\Sh{\mathop{\mathit{Sh}}\nolimits}
\def\NL{\mathop{N\!L}\nolimits}
\def\CH{\mathop{\mathrm{CH}}\nolimits}
\def\proetale{{pro\text{-}\acute{e}tale}}
\def\etale{{\acute{e}tale}}
\def\QCoh{\mathit{QCoh}}
\def\Ker{\mathop{\mathrm{Ker}}}
\def\Im{\mathop{\mathrm{Im}}}
\def\Coker{\mathop{\mathrm{Coker}}}
\def\Coim{\mathop{\mathrm{Coim}}}

% Boxtimes
%
\DeclareMathSymbol{\boxtimes}{\mathbin}{AMSa}{"02}

%
% Macros for moduli stacks/spaces
%
\def\QCohstack{\mathcal{QC}\!\mathit{oh}}
\def\Cohstack{\mathcal{C}\!\mathit{oh}}
\def\Spacesstack{\mathcal{S}\!\mathit{paces}}
\def\Quotfunctor{\mathrm{Quot}}
\def\Hilbfunctor{\mathrm{Hilb}}
\def\Curvesstack{\mathcal{C}\!\mathit{urves}}
\def\Polarizedstack{\mathcal{P}\!\mathit{olarized}}
\def\Complexesstack{\mathcal{C}\!\mathit{omplexes}}
% \Pic is the operator that assigns to X its picard group, usage \Pic(X)
% \Picardstack_{X/B} denotes the Picard stack of X over B
% \Picardfunctor_{X/B} denotes the Picard functor of X over B
\def\Pic{\mathop{\mathrm{Pic}}\nolimits}
\def\Picardstack{\mathcal{P}\!\mathit{ic}}
\def\Picardfunctor{\mathrm{Pic}}
\def\Deformationcategory{\mathcal{D}\!\mathit{ef}}
\usepackage{cleveref}
\crefname{originalequation}{original equation}{original equations}
\Crefname{originalequation}{Original equation}{Original equations}
\makeatletter
\newlabel{equation-bijective-orbits}{{06T6}{source}{}{sourcepointer-1}{}}
\newlabel{equation-bijective-orbits@cref}{{[originalequation][0][]06T6}{[1][0][]source}{}{}{}}
\makeatother
\crefname{originalequation}{original equation}{original equations}
\Crefname{originalequation}{Original equation}{Original equations}
\makeatletter
\newlabel{equation-sequence}{{06S6}{source}{}{sourcepointer-2}{}}
\newlabel{equation-sequence@cref}{{[originalequation][0][]06S6}{[1][0][]source}{}{}{}}
\makeatother
\crefname{originalreference}{original reference}{original references}
\Crefname{originalreference}{Original reference}{Original references}
\makeatletter
\newlabel{lemma-action-linear}{{06SU}{source}{}{sourcepointer-3}{}}
\newlabel{lemma-action-linear@cref}{{[originalreference][0][]06SU}{[1][0][]source}{}{}{}}
\makeatother
\crefname{originalreference}{original reference}{original references}
\Crefname{originalreference}{Original reference}{Original references}
\makeatletter
\newlabel{lemma-essential-surjection}{{06H0}{source}{}{sourcepointer-4}{}}
\newlabel{lemma-essential-surjection@cref}{{[originalreference][0][]06H0}{[1][0][]source}{}{}{}}
\makeatother
\crefname{originalreference}{original reference}{original references}
\Crefname{originalreference}{Original reference}{Original references}
\makeatletter
\newlabel{lemma-k-linear-differential}{{06IJ}{source}{}{sourcepointer-5}{}}
\newlabel{lemma-k-linear-differential@cref}{{[originalreference][0][]06IJ}{[1][0][]source}{}{}{}}
\makeatother
\crefname{originalreference}{original reference}{original references}
\Crefname{originalreference}{Original reference}{Original references}
\makeatletter
\newlabel{lemma-lifting-section}{{06IS}{source}{}{sourcepointer-6}{}}
\newlabel{lemma-lifting-section@cref}{{[originalreference][0][]06IS}{[1][0][]source}{}{}{}}
\makeatother
\crefname{originalreference}{original reference}{original references}
\Crefname{originalreference}{Original reference}{Original references}
\makeatletter
\newlabel{lemma-minimal-versal}{{06T5}{source}{}{sourcepointer-7}{}}
\newlabel{lemma-minimal-versal@cref}{{[originalreference][0][]06T5}{[1][0][]source}{}{}{}}
\makeatother
\crefname{originalreference}{original reference}{original references}
\Crefname{originalreference}{Original reference}{Original references}
\makeatletter
\newlabel{lemma-smallest-where-descends-versal}{{06T2}{source}{}{sourcepointer-8}{}}
\newlabel{lemma-smallest-where-descends-versal@cref}{{[originalreference][0][]06T2}{[1][0][]source}{}{}{}}
\makeatother
\crefname{originalreference}{original reference}{original references}
\Crefname{originalreference}{Original reference}{Original references}
\makeatletter
\newlabel{lemma-tangent-space-vector-space}{{06IH}{source}{}{sourcepointer-9}{}}
\newlabel{lemma-tangent-space-vector-space@cref}{{[originalreference][0][]06IH}{[1][0][]source}{}{}{}}
\makeatother
\title{The minimal versal formal object of a predeformation category satisfying (S2) has bijective orbits}
\author{The Stacks Project Authors}
\date{stacks-project:formal-defos.tex:local-snapshot-20261012; GNU FDL 1.2 or later}
\begin{document}
\maketitle
\noindent\textit{Editorial note.} All delivered mathematical text is the author's own, copied from the pinned Stacks Project chapter file formal-defos.tex; no statement or proof step was rewritten, reordered or removed.\par\medskip
\noindent\textit{Editorial note.} The counted proof is the author's complete argument: the diagram interpretation through the vector space structure, the k-linearity of the differential and the action by addition, the construction of the subring on which the derivation vanishes with its cotangent space computation, and the final contradiction with minimality via the lifting lemma.\par\medskip
\noindent\textit{Editorial note.} The label lemma-construct-bijective-orbits names the counted statement, and the proof environment immediately following it in the source is delivered verbatim as its proof.\par\medskip
\noindent\textit{Editorial note.} The counted statement is carried by the author's own lemma environment in the chapter source.\par\medskip
\noindent\textit{Editorial note.} Pre-open gate: python operations/agent-authoring/tools/precard\_gate.py --source formal-defos.tex --label lemma-construct-bijective-orbits --start 4743 --end 4834 --local-id 28012 returned RESULT=OK (L28012 free, occupancy FREE, public interval clear). Candidate taken from sources/stacks-formal-defos/candidates-halfA.json (second exclusive pool assigned by the Lead on 2026-10-13).\par\medskip
\noindent\textit{Editorial note.} The statement and proof are copied from the pinned chapter file; the cited inputs are kept as external pointers with their canonical Stacks tags.\par\medskip
\noindent\textit{Allowed standard prerequisites (editorial).}\begin{itemize}
\item Existence and minimality of a versal formal object. Same chapter, the lemmas producing a minimal versal formal object of the predeformation category. Supplies the minimal versal object xi over R used throughout the proof and the minimality property invoked at the end.
\item Vector space structure and linearity of the differential. Same chapter, the lemma giving the tangent space its k-vector space structure, the lemma that the differential of a formal object is k-linear, and the lemma identifying the action of the derivations of k with addition. Supply the interpretation of the displayed diagram: the kernel K, the linearity of the differential and the action by addition.
\item Lifting of a section along a morphism hitting the trivial deformation. Same chapter, the lemma allowing a morphism x -> y to be produced once the pushforward of the trivial object is reached. Supplies the morphism x -> y obtained from the triviality of x\_epsilon.
\item Descending the versal object along a nonsurjective ring map. Same chapter, part (1) of the lemma characterizing the smallest ring over which a versal object descends. Supplies the contradiction with minimality once the ring map R'/m\_R\textasciicircum{}2 -> R/m\_R\textasciicircum{}2 is seen to be nonsurjective.
\end{itemize}\medskip
\ifdefined\hypertarget\hypertarget{sourcepointer-1}{}\fi
\noindent\textit{Source reference.} Original equation 06T6: sources/stacks-formal-defos/pinned-original-source.tex, source lines 4623--4623. Prerequisite pointer invoked at proof line(s) 4747; it names a standard fact of the same chapter that the counted argument uses, and every step of the counted proof built on it is retained verbatim in the delivered body.\par\medskip
\ifdefined\hypertarget\hypertarget{sourcepointer-2}{}\fi
\noindent\textit{Source reference.} Original equation 06S6: sources/stacks-formal-defos/pinned-original-source.tex, source lines 808--808. Prerequisite pointer invoked at proof line(s) 4775; it names a standard fact of the same chapter that the counted argument uses, and every step of the counted proof built on it is retained verbatim in the delivered body.\par\medskip
\ifdefined\hypertarget\hypertarget{sourcepointer-3}{}\fi
\noindent\textit{Source reference.} Original reference 06SU: sources/stacks-formal-defos/pinned-original-source.tex, source lines 3811--3811. Prerequisite pointer invoked at proof line(s) 4763; it names a standard fact of the same chapter that the counted argument uses, and every step of the counted proof built on it is retained verbatim in the delivered body.\par\medskip
\ifdefined\hypertarget\hypertarget{sourcepointer-4}{}\fi
\noindent\textit{Source reference.} Original reference 06H0: sources/stacks-formal-defos/pinned-original-source.tex, source lines 887--887. Prerequisite pointer invoked at proof line(s) 4782; it names a standard fact of the same chapter that the counted argument uses, and every step of the counted proof built on it is retained verbatim in the delivered body.\par\medskip
\ifdefined\hypertarget\hypertarget{sourcepointer-5}{}\fi
\noindent\textit{Source reference.} Original reference 06IJ: sources/stacks-formal-defos/pinned-original-source.tex, source lines 3716--3716. Prerequisite pointer invoked at proof line(s) 4760; it names a standard fact of the same chapter that the counted argument uses, and every step of the counted proof built on it is retained verbatim in the delivered body.\par\medskip
\ifdefined\hypertarget\hypertarget{sourcepointer-6}{}\fi
\noindent\textit{Source reference.} Original reference 06IS: sources/stacks-formal-defos/pinned-original-source.tex, source lines 3212--3212. Prerequisite pointer invoked at proof line(s) 4824; it names a standard fact of the same chapter that the counted argument uses, and every step of the counted proof built on it is retained verbatim in the delivered body.\par\medskip
\ifdefined\hypertarget\hypertarget{sourcepointer-7}{}\fi
\noindent\textit{Source reference.} Original reference 06T5: sources/stacks-formal-defos/pinned-original-source.tex, source lines 4558--4558. Prerequisite pointer invoked at proof line(s) 4752; it names a standard fact of the same chapter that the counted argument uses, and every step of the counted proof built on it is retained verbatim in the delivered body.\par\medskip
\ifdefined\hypertarget\hypertarget{sourcepointer-8}{}\fi
\noindent\textit{Source reference.} Original reference 06T2: sources/stacks-formal-defos/pinned-original-source.tex, source lines 4375--4375. Prerequisite pointer invoked at proof line(s) 4831; it names a standard fact of the same chapter that the counted argument uses, and every step of the counted proof built on it is retained verbatim in the delivered body.\par\medskip
\ifdefined\hypertarget\hypertarget{sourcepointer-9}{}\fi
\noindent\textit{Source reference.} Original reference 06IH: sources/stacks-formal-defos/pinned-original-source.tex, source lines 3669--3669. Prerequisite pointer invoked at proof line(s) 4757; it names a standard fact of the same chapter that the counted argument uses, and every step of the counted proof built on it is retained verbatim in the delivered body.\par\medskip
\begin{lemma}
\label{lemma-construct-bijective-orbits}
Let $\mathcal{F}$ be a predeformation category satisfying
(S2) which has a versal formal object. Then its minimal versal
formal object satisfies (\ref{equation-bijective-orbits}).
\end{lemma}

\begin{proof}
Let $\xi$ be a minimal versal formal object for $\mathcal{F}$, see
Lemma \ref{lemma-minimal-versal}.
Say $\xi$ lies over $R \in \Ob(\widehat{\mathcal{C}}_\Lambda)$.
In order to parse (\ref{equation-bijective-orbits}) we point out
that $T\mathcal{F}$ has a natural $k$-vector space structure
(see
Lemma \ref{lemma-tangent-space-vector-space}),
that $d\underline{\xi} : \text{Der}_\Lambda(R, k) \to T\mathcal{F}$
is linear (see
Lemma \ref{lemma-k-linear-differential}),
and that the action of $\text{Der}_\Lambda(k, k)$ is
given by addition (see
Lemma \ref{lemma-action-linear}).
Consider the diagram
$$
\xymatrix{
& \Hom_k(\mathfrak m_R/\mathfrak m_R^2, k) \\
K \ar[r] & \text{Der}_\Lambda(R, k) \ar[r]^{d\underline{\xi}} \ar[u] &
T\mathcal{F} \\
& \text{Der}_\Lambda(k, k) \ar[u] \ar[ru]
}
$$
The vector space $K$ is the kernel of $d\underline{\xi}$.
Note that the middle column is exact in the middle as it is dual to the
sequence (\ref{equation-sequence}). If (\ref{equation-bijective-orbits})
fails, then we can find a nonzero element $D \in K$ which
does not map to zero in $\Hom_k(\mathfrak m_R/\mathfrak m_R^2, k)$.
This means there exists an $t \in \mathfrak m_R$ such that
$D(t) = 1$. Set $R' = \{a \in R \mid D(a) = 0\}$. As $D$ is a derivation
this is a subring of $R$. Since $D(t) = 1$ we see that $R' \to k$
is surjective (compare with the proof of
Lemma \ref{lemma-essential-surjection}).
Note that $\mathfrak m_{R'} = \Ker(D : \mathfrak m_R \to k)$
is an ideal of $R$ and $\mathfrak m_R^2 \subset \mathfrak m_{R'}$. Hence
$$
\mathfrak m_R/\mathfrak m_R^2 =
\mathfrak m_{R'}/\mathfrak m_R^2 + k\overline{t}
$$
which implies that the map
$$
R'/\mathfrak m_R^2 \times_k k[\epsilon] \to R/\mathfrak m_R^2
$$
sending $\epsilon$ to $\overline{t}$ is an isomorphism. In particular
there is a map $R/\mathfrak m_R^2 \to R'/\mathfrak m_R^2$.

\medskip\noindent
Let $\xi \to y$ be a morphism lying over $R \to R/\mathfrak m_R^2$.
Let $y \to x$ be a morphism lying over
$R/\mathfrak m_R^2 \to R'/\mathfrak m_R^2$.
Let $y \to x_\epsilon$ be a morphism lying over
$R/\mathfrak m_R^2 \to k[\epsilon]$. Let $x_0$ be the unique (up to unique
isomorphism) object of $\mathcal{F}$ over $k$.
By the axioms of a category cofibred in groupoids we obtain
a commutative diagram
$$
\vcenter{
\xymatrix{
y \ar[r] \ar[d] & x_\epsilon \ar[d] \\
x \ar[r]   & x_0
}
}
\quad\text{lying over}\quad
\vcenter{
\xymatrix{
R'/\mathfrak m_R^2 \times_k k[\epsilon] \ar[r] \ar[d] & k[\epsilon] \ar[d] \\
R'/\mathfrak m_R^2 \ar[r] & k.
}
}
$$
Because $D \in K$ we see that $x_\epsilon$ is isomorphic
to $0 \in \mathcal{F}(k[\epsilon])$, i.e., $x_\epsilon$ is
the pushforward of $x_0$ via $k \to k[\epsilon], a \mapsto a$.
Hence by
Lemma \ref{lemma-lifting-section}
we see that there exists a morphism $x \to y$. Since
$\text{length}_\Lambda(R'/\mathfrak m_R^2) <
\text{length}_\Lambda(R/\mathfrak m_R^2)$
the corresponding ring map $R'/\mathfrak m_R^2 \to R/\mathfrak m_R^2$
is not surjective. This contradicts the minimality of
$\xi/R$, see part (1) of
Lemma \ref{lemma-smallest-where-descends-versal}.
This contradiction shows that such a $D$ cannot exist, hence
we win.
\end{proof}

\end{document}
